%-----------------------------------------------------------------------------------------------------------------------------------------------%
%	The MIT License (MIT)
%
%	Copyright (c) 2015 Jan Küster
%
%	Permission is hereby granted, free of charge, to any person obtaining a copy
%	of this software and associated documentation files (the "Software"), to deal
%	in the Software without restriction, including without limitation the rights
%	to use, copy, modify, merge, publish, distribute, sublicense, and/or sell
%	copies of the Software, and to permit persons to whom the Software is
%	furnished to do so, subject to the following conditions:
%	
%	THE SOFTWARE IS PROVIDED "AS IS", WITHOUT WARRANTY OF ANY KIND, EXPRESS OR
%	IMPLIED, INCLUDING BUT NOT LIMITED TO THE WARRANTIES OF MERCHANTABILITY,
%	FITNESS FOR A PARTICULAR PURPOSE AND NONINFRINGEMENT. IN NO EVENT SHALL THE
%	AUTHORS OR COPYRIGHT HOLDERS BE LIABLE FOR ANY CLAIM, DAMAGES OR OTHER
%	LIABILITY, WHETHER IN AN ACTION OF CONTRACT, TORT OR OTHERWISE, ARISING FROM,
%	OUT OF OR IN CONNECTION WITH THE SOFTWARE OR THE USE OR OTHER DEALINGS IN
%	THE SOFTWARE.
%	
%
%-----------------------------------------------------------------------------------------------------------------------------------------------%


%============================================================================%
%
%	DOCUMENT DEFINITION
%
%============================================================================%

%we use article class because we want to fully customize the page and dont use a cv template
\documentclass[10pt,A4]{article}	


%----------------------------------------------------------------------------------------
%	ENCODING
%----------------------------------------------------------------------------------------

%we use utf8 since we want to build from any machine
\usepackage[utf8]{inputenc}		

%----------------------------------------------------------------------------------------
%	LOGIC
%----------------------------------------------------------------------------------------

% provides \isempty test
\usepackage{xifthen}

%----------------------------------------------------------------------------------------
%	FONT
%----------------------------------------------------------------------------------------

% some tex-live fonts - choose your own

%\usepackage[defaultsans]{droidsans}
%\usepackage[default]{comfortaa}
%\usepackage{cmbright}
%\usepackage[default]{raleway}
%\usepackage{fetamont}
%\usepackage[default]{gillius}
%\usepackage[light,math]{iwona}
%\usepackage[thin]{roboto} 

% set font default
\renewcommand*\familydefault{\sfdefault} 	
\usepackage[T1]{fontenc}

% more font size definitions
\usepackage{moresize}		


%----------------------------------------------------------------------------------------
%	PAGE LAYOUT  DEFINITIONS
%----------------------------------------------------------------------------------------

%debug page outer frames
%\usepackage{showframe}			


%define page styles using geometry
\usepackage[a4paper]{geometry}		

% for example, change the margins to 2 inches all round
\geometry{top=20mm, bottom=20mm, left=1.5cm, right=1.5cm} 	

%use customized header
\usepackage{fancyhdr}
\pagestyle{fancy}
\fancyhf{}

%use our custom fancy header definitions
\fancyfoot[R]{\thepage}
\renewcommand{\headrulewidth}{0pt}
\renewcommand{\footrulewidth}{0pt}

%less space between header and content
\setlength{\headheight}{-5pt}		


%customize entries left, center and right
%\lhead{}
%\chead{}
%\rhead{}


%indentation is zero
\setlength{\parindent}{0mm}

%----------------------------------------------------------------------------------------
%	TABLE /ARRAY DEFINITIONS
%---------------------------------------------------------------------------------------- 

%for layouting tables
\usepackage{multicol}			
\usepackage{multirow}

%extended aligning of tabular cells
\usepackage{array}

\newcolumntype{x}[1]{%
>{\raggedleft\hspace{0pt}}p{#1}}%


%----------------------------------------------------------------------------------------
%	GRAPHICS DEFINITIONS
%---------------------------------------------------------------------------------------- 

%for header image
\usepackage{graphicx}

%for floating figures
\usepackage{wrapfig}
\usepackage{float}
%\floatstyle{boxed} 
%\restylefloat{figure}

%for drawing graphics		
\usepackage{tikz}				
\usetikzlibrary{shapes, backgrounds,mindmap, trees}


%----------------------------------------------------------------------------------------
%	Color DEFINITIONS
%---------------------------------------------------------------------------------------- 

%accent color
\definecolor{sectcol}{RGB}{252,173,98}

\definecolor{sectcoldark}{RGB}{125, 62, 3}

%dark background color
\definecolor{bgcol}{RGB}{110,110,110}

%light background / accent color
\definecolor{softcol}{RGB}{225,225,225}

% Links
\usepackage{color}
\usepackage{hyperref}
\hypersetup{
    colorlinks=true,
    urlcolor=sectcoldark
    }



%============================================================================%
%
%
%	DEFINITIONS
%
%
%============================================================================%

%----------------------------------------------------------------------------------------
% 	HEADER
%----------------------------------------------------------------------------------------

%remove pagenum
% \renewcommand{\thepage}{}	

%remove section num		
\renewcommand{\thesection}{}			

%----------------------------------------------------------------------------------------
% 	ARROW GRAPHICS in Tikz
%----------------------------------------------------------------------------------------

% a six pointed arrow poiting to the left
\newcommand{\tzlarrow}{(0,0) -- (0.2,0) -- (0.3,0.2) -- (0.2,0.4) -- (0,0.4) -- (0.1,0.2) -- cycle;}	

% include the left arrow into a tikz picture
% param1: fill color
%
\newcommand{\larrow}[1]
{\begin{tikzpicture}[scale=0.58]
	 \filldraw[fill=#1!100,draw=#1!100!black]  \tzlarrow
 \end{tikzpicture}
}

% a six pointed arrow poiting to the right
\newcommand{\tzrarrow}{ (0,0.2) -- (0.1,0) -- (0.3,0) -- (0.2,0.2) -- (0.3,0.4) -- (0.1,0.4) -- cycle;}

% include the right arrow into a tikz picture
% param1: fill color
%
\newcommand{\rarrow}
{\begin{tikzpicture}[scale=0.7]
	\filldraw[fill=sectcol!100,draw=sectcol!100!black] \tzrarrow
 \end{tikzpicture}
}



%----------------------------------------------------------------------------------------
%	custom sections
%----------------------------------------------------------------------------------------

% create a coloured box with arrow and title as cv section headline
% param 1: section title
%
\newcommand{\cvsection}[1]
{
	\begin{center}
		\large\textcolor{sectcol}{\textbf{#1}}
	\end{center}
}

%create a coloured arrow with title as cv meta section section
% param 1: meta section title
%
\newcommand{\metasection}[2]
{
%\begin{tabular*}{1\textwidth}{r r}
\footnotesize{#2} \hspace*{\fill} \footnotesize{#1}\\[1pt]
%\end{tabular*}
}

%----------------------------------------------------------------------------------------
%	 CV EVENT
%----------------------------------------------------------------------------------------

% creates a stretched box as cv entry headline followed by some paragraphs about 
% the work you did
% param 1:	event time i.e. 2014 or 2011-2014 etc.
% param 2:	event name (what did you do?)
% param 3:	institution (where did you work / study)
% param 4:	list of paragraphs
%
\newcommand{\cvevent}[4]
{

\begin{tabular*}{1\textwidth}{p{13.6cm}  x{3.9cm}}
	\textbf{#2} - \textcolor{bgcol}{#3} &   \vspace{2.5pt}\textcolor{sectcol}{#1}
\end{tabular*}

\vspace{-8pt}
\textcolor{softcol}{\hrule}
\vspace{6pt}

	\foreach \desc in {#4}{
		$\cdot$ \desc\\[3pt]
	}
	
\vspace{3pt}
}

% creates a stretched box as 
\newcommand{\cveventmeta}[2]
{
	\mbox{\mystrut \hspace{87pt}\textit{#1}}\\
	#2
}

%----------------------------------------------------------------------------------------
% CUSTOM STRUT FOR EMPTY BOXES
%----------------------------------------- -----------------------------------------------
\newcommand{\mystrut}{\rule[-.3\baselineskip]{0pt}{\baselineskip}}

%============================================================================%
%
%
%
%	DOCUMENT CONTENT
%
%
%
%============================================================================%
\begin{document}



%----------------------------------------------------------------------------------------
%	HEADER IMAGE
%----------------------------------------------------------------------------------------

%\begin{figure}[H]
%\begin{flushright}
%	\includegraphics[trim= 320 130 460 210,clip,width=0.2\linewidth]{myfoto.jpg}	%trimming relative to image size!
%\end{flushright}
%\end{figure}


%---------------------------------------------------------------------------------------
%	TITLE HEADLINE
%----------------------------------------------------------------------------------------
\vspace{-8pt}
\begin{center}
	\HUGE \textsc{Raul Otaño Hurtado} \textcolor{sectcol}{\rule[-1mm]{1mm}{0.9cm}} \textsc{Resume}\\[2pt]
	\small Software Engineer
\end{center}



\vspace{6pt}


%---------------------------------------------------------------------------------------
%	META SECTION
%----------------------------------------------------------------------------------------
\metasection{Madrid, Spain}{\textbf{Status:} B.Sc. Computer Science}
\metasection{\href{https://rulyotano.com}{rulyotano.com}}{\textbf{Fields:} Software Engineering} 
\metasection{<<emailraul>>}{\textbf{Languages:} English (Professional), Spanish (Native), Portuguese (Professional)}
\metasection{<<phoneraul>>}{\textbf{Activities:} Swimming, Running, Open Source, Fitness, Babysitting}
\metasection{}{\textbf{Social:} \href{https://github.com/rulyotano}{Github}, \href{https://www.linkedin.com/in/raulotanohurtado}{Linkedin}, \href{https://stackoverflow.com/users/1655482}{Stackoverflow}, \href{https://www.hackerrank.com/rulyotano}{Hackerrank}, \href{https://medium.com/@rulyotano}{Medium}}
\vspace{-2pt}
\textcolor{softcol}{\hrule}
\vspace{6pt}

\normalsize

%---------------------------------------------------------------------------------------
%	SUMMARY (optional)
%----------------------------------------------------------------------------------------
\vspace{-6pt}
\cvsection{Summary}
I'm a passionate software engineer who likes clean code, good practices, and SOLID principles. I love coding, algorithms, data structures, and learning about technology.\\

I like system design topics and would love to help to create scalable and efficient solutions. I'm also enthusiastic about contributing to open-source projects, enhancing libraries, and assisting other developers.\\

I would like to collaborate and create outstanding software with you!\\


%============================================================================%
%
%	CV SECTIONS AND EVENTS (MAIN CONTENT)
%
%============================================================================%

%---------------------------------------------------------------------------------------
%	EXPERIENCE
%----------------------------------------------------------------------------------------
\cvsection{Experience}

\cvevent{May 2025 - present}{Senior Software Engineer}{Payoneer (\href{https://www.payoneer.com}{payoneer.com})}{
	{Design and implement new solutions for new requirements or integrations with external service providers}
    {Update our codebase from legacy to a more standardized architecture. E.g., make services compatible with Docker, stateless, and integration with k8s}
}

\cvevent{March 2021 - May 2025}{Software Engineer II (Backed .NET)}{Docplanner (\href{https://www.docplanner.com/about-us}{docplanner.com})}{
	{Design, develop, test, and monitor new features to help doctors manage their patients and bookings.},
	{Enhance the current code base by improving the performance, creating new tests, or doing meaningful refactoring to the code.},
	{Keep the lights on by investigating and finding the root cause of problems in our platform and ensuring a fast issue resolution.},
	{Had the opportunity to participate in the designing, coding, deployment, monitoring, and support of new cross-company services used by other teams.}, 
	{Be supported and try to help everyone as much as I can by sharing my domain knowledge and experience with the final goal of helping the company and consequently its clients.}
}

\cvevent{July 2020 - March 2021}{Senior Software Engineer for Openbank}{Mimacom Iberica (\href{https://mimacom.com/es}{mimacom.com})}{
	{I served as a Senior Software Developer at Openbank, a client of Mimacom. My role primarily focused on front-end development using React.js. The primary challenge was to develop a new onboarding application while adhering to best practice standards and incorporating Openbank technologies. A portion of the application required integration into the mobile app, which presented an additional challenge. I also conducted numerous job interviews for candidates interested in joining Mimacom.}
}

\cvevent{March 2018 - June 2020}{Fullstack Software Engineer}{Consolfreight (\href{https://consolfreight.com/}{consolfreight.com})}{
	{I served as a full-stack engineer to successfully complete and launch the ConsolFreight web application. The primary objective of this system was to facilitate the sale of cargo routes to individuals and businesses looking to ship their goods. One of the major challenges involved translating complex industry requirements into a coherent domain definition that could be implemented in the application. I utilized a combination of technologies, including Aspnet MVC, MVC Core, ReactJs/Redux, and also played a key role in setting up CI/CD pipelines using Azure.}
}

\cvevent{September 2016 - 2018}{Lead Frontend Software Engineer}{Solinftec (\href{https://solinftec.com/}{solinftec.com})}{
	{Developed the web frontend system to display agriculture field data to users. The system displays maps showing the position and status of equipment. It also provides real-time weather condition information. The system sends alerts to users when specific parameters deviate from the expected range. Users can filter and manipulate the data. One of the key challenges was integrating various technologies, such as AngularJS and Google Maps. Additionally, we successfully migrated from Google Maps to Mapbox and from AngularJS to ReactJS over time.}
}

\cvevent{September 2014 - 2016}{Fullstack Software Engineer}{Xmovo (\href{https://www.linkedin.com/company/xmovo-group-inc/about}{xmovo.com})}{
	{I collaborated with the Xmovo team to develop a new social network website focused on training and fitness. This ambitious project featured a range of components, including a feeds section, notification system, chat functionality, e-commerce for training programs and memberships, live video training classes, and more. Unfortunately, the project was not successful, and the owners decided to close it.},
	{Key Technologies: C\#, ASP.NET MVC 5, Entity Framework 6, AngularJS, React.js, SignalR, Domain-Driven Design (DDD), SOLID principles, and more.}
}

\cvevent{2010 - 2013}{Freelancer and Previous jobs}{Many companies (\href{https://www.synnotech.de/}{Synotech AG}, doiTTeam, \href{https://www.linkedin.com/company/vworker}{vWorker}, ...)}{
	{During my university years, I primarily worked as a freelance developer for various companies.},
	{I developed numerous WPF applications and gained proficiency in best practices and design patterns, with MVVM being a prominent example, particularly in the context of WPF and Silverlight.},
	{I collaborated with colleagues to establish the doiTTeam, an initiative focused on organizing our freelance work.},
	{I successfully completed a couple of projects during my internship with a specific company I was assigned to.},
	{During this period, I acquired valuable skills in real-world application development, primarily focusing on .NET technologies, including C\#, WPF, Windows Forms, NHibernate, ActiveRecord, Entity Framework, various ORMs, and database systems like SQL Server and SQLite.}
}

%---------------------------------------------------------------------------------------
%	EDUCATION SECTION
%--------------------------------------------------------------------------------------
\cvsection{Education}

\cvevent{2008 - 2013}{Graduated as B.Sc. Computer Science}{\href{https://matcom.uh.cu}{Havana University}}{
	{I graduated with a strong foundation in Computer Science, where I gained knowledge in a wide range of subjects.},
	{My coursework covered fundamental topics such as mathematics, logic, and algebra, providing a solid grounding for my studies.},
	{In addition, I dug into various Computer Science-related subjects, including programming, compilers, artificial intelligence, databases, networking, optimizations, and more.}
}

\cvevent{2020 - present}{Licenses and Certifications}{Udemy, Linkedin, HackerRank}{
	{Docker Mastery: with Kubernetes +Swarm from a Docker Captain (\href{https://ude.my/UC-1cd13819-15b9-497f-b64a-bdcf46a5486b}{Udemy})},
	{Software Architecture Foundations (\href{https://www.linkedin.com/learning/certificates/bd4c1327e9f788cc7965dcded5641e21e2fbc9855374203a2f84ad7bb7910fab}{Linkedin})},
	{Software Architecture: Domain-Driven Design (\href{https://www.linkedin.com/learning/certificates/e337a7bc24c5ab4bb34600bed69510b387601cc8c75df9f43c73bde385a4ef67}{Linkedin})},
	{Microservices Architecture and Implementation on .NET 5 Course (\href{https://ude.my/UC-1a9390d9-ffee-4789-9c27-6bdbb35fa1f5}{Udemy})},
	{Problem Solving (Intermediate) Certificate (\href{https://www.hackerrank.com/certificates/5b4c0d75cbc9}{HackerRank})}
}


%---------------------------------------------------------------------------------------
%	OPEN SOURCE SECTION
%--------------------------------------------------------------------------------------
\cvsection{Open Source}

\cvevent{}{Tree Extended}{\href{https://github.com/rulyotano/tree-extended}{rulyotano/tree-extended}}{
	{Tree Extended is an NPM package that can be installed globally, serving as a command-line tool for your terminal. It works seamlessly on Linux, Windows, and macOS. This tool is designed to generate text directory trees, similar to the "tree" command. What sets Tree Extended apart are its customizability features. You can tailor the output by applying filters, limiting the depth level of the tree, and choosing from a variety of character sets, including UTF-8, ASCII, or UTF-8 with icons. The project follows an automated deployment process to NPM using GitHub pipelines, ensuring accessibility to the wider developer community.}
}

\cvevent{}{Tree Extended VsCode Plugin}{\href{https://github.com/rulyotano/tree-extended-vscode}{rulyotano/tree-extended-vscode}}{
	{This one is a Visual Studio Code (VSCode) extension that leverages the capabilities of the Tree Extended library to generate directory trees directly within the VSCode environment. This plugin streamlines the process, making it more user-friendly and accessible to everyone using VSCode for their development work. The project adopts an automated deployment process to the Visual Studio Code Marketplace, utilizing GitHub pipelines for seamless and efficient distribution.}
}

\cvevent{}{Rulyotano.Crosscuting}{\href{https://github.com/rulyotano/Rulyotano.CrossCutting}{Rulyotano.CrossCutting}, \href{https://github.com/rulyotano/rulyotano.crosscutting.js}{rulyotano.crosscutting.js}}{
	{These crosscutting libraries are my personal implementations, created during my available time. They automate the generation of NuGet packages for .NET and NPM packages for JavaScript. Within these libraries, you'll find a range of resources, including common algorithms and data structures that can be useful in various software projects. The most interesting one, in my opinion, is the \href{https://rulyotano.com/demos/bezier}{interpolation tool}, implemented for both, dotnet and js. We can get a Bezier curve that passes through a set of points previously defined. Also, there is a \href{https://github.com/rulyotano/wpf-bezier-interpolation}{WPF sample application}.}
}

\cvevent{}{Minesweeper}{\href{https://github.com/rulyotano/minesweeper}{rulyotano/minesweeper}, \href{https://minesweeper.rulyotano.com}{minesweeper.rulyotano.com}}{
	{This project began as a personal endeavor to create a Minesweeper game using the React.js and Redux patterns. What initially started as a simple example evolved into a fully working application. To put architectural knowledge to the test, I took on the challenge of dockerizing the application and expanded it by incorporating additional services, including a Redis database and a Traefik reverse proxy.},
        {The entire project was deployed on a Docker Swarm, hosted on a cost-effective VPS. Leveraging Docker Swarm ensures scalability and resilience for the application.},
	{The project remains a playground for continuous improvement. I actively use and enhance it, incorporating new features and experimenting with various aspects of software development.}
}

\cvevent{}{Httpd-spa Docker Image}{\href{https://github.com/rulyotano/httpd-spa}{rulyotano/httpd-spa}}{
	{This project provides a straightforward yet highly useful Docker image configuration. It's pre-configured to function as an Apache server specialized in serving single-page applications (SPAs). The specific challenge this image addresses is common in the context of SPAs: when SPAs manage routing internally, and the app is reloaded in sub-routes, the default httpd image attempts to locate other resources.},
        {Purpose: The image's purpose is to simplify the deployment of SPAs and ensure smooth routing, regardless of the internal management of routes within the SPA.},
	{Ease of Use: By using this Docker image, developers can quickly set up an Apache server configured to serve SPAs without the complications of managing route-related issues}
}

\cvevent{}{Cv}{\href{https://github.com/rulyotano/cv}{rulyotano/cv}}{
	{This open-source project offers a LaTeX-based CV template that is designed for professional use. The template is intended to serve as a foundation that can be adapted to various other LaTeX templates, making it highly versatile for individuals in need of a well-structured and customizable CV.},
        {Current State: The project is a continuous work in progress, reflecting the evolving needs of users and the incorporation of new design elements and features to enhance the template's usability.},
	{Versatility: The LaTeX CV template project aims to provide a flexible and adaptable solution for individuals looking to create a visually appealing and professionally formatted CV.}
}

\cvevent{}{Mui-profile}{\href{https://github.com/rulyotano/mui-profile}{rulyotano/mui-profile}}{
	{My personal profile website is built using React.js and Nuxt.js, showcasing a dynamic and interactive representation of my professional identity. Hosted on GitHub Pages, it provides a platform for sharing my skills, projects, and achievements, offering a comprehensive overview of my capabilities and experiences.},
        {This project is designed to be versatile and adaptable, making it a valuable resource for anyone looking to create their own profile website. The website's codebase can be reused and customized for various purposes.},
	{I have authored \href{https://betterprogramming.pub/get-your-personal-website-for-free-create-it-with-reactjs-b7e3c3c874b4}{an informative article} detailing the development and setup of this website, which serves as a helpful guide for others interested in creating their own online profile.}
}

\cvevent{}{Coding interview preparation checklist}{\href{https://github.com/rulyotano/coding-interview-preparation-checklist}{rulyotano/coding-interview-preparation-checklist}}{
	{The "Coding Interview Preparation Checklist" is an open-source project primarily comprising a curated collection of references to valuable materials and resources for coding interview preparation. It serves as a comprehensive guide that directs individuals to a wide range of topics, strategies, and techniques essential for excelling in technical interviews.}
}

\cvevent{}{Other Repositories}{\href{https://github.com/rulyotano?tab=repositories}{rulyotano/repositories}}{
	{My GitHub profile offers a diverse collection of open-source projects and repositories, covering a wide range of topics and technologies. You are welcome to explore my repositories to discover various interesting projects, code samples, and contributions.}
}

%============================================================================%
%
%
%
%	DOCUMENT END
%
%
%
%============================================================================%
\end{document}
